% Generated by code/fill_numbers.py from chuong3_end.tpl -- edit the template, not this file.
% ----------------------------------------------------------------------------
\section{Thảo luận}\label{sec:discussion}

\paragraph{Trả lời các câu hỏi nghiên cứu.} Các thực nghiệm của chương cho phép trả lời bốn câu hỏi nêu ở Mở đầu.

\emph{(CH1) Độ chính xác.} Trên dữ liệu hai chiều, PWL-SVM dạng HH đạt độ chính xác gần SVM hạt nhân Gauss, trừ trên bàn cờ (\cref{tab:exp2d}). Trên mười bộ dữ liệu chuẩn, PWL-C-SVM cộng tính đứng giữa SVM tuyến tính và SVM hạt nhân Gauss: nó kém SVM hạt nhân Gauss trung bình 1{,}4 điểm phần trăm, một khác biệt có ý nghĩa thống kê, và hơn SVM tuyến tính trung bình 1{,}7 điểm, với phần cải thiện tập trung ở những bộ dữ liệu phi tuyến rõ rệt (\cref{tab:bench,tab:wilcoxon}). Nhận định của bài báo gốc rằng PWL-SVM chính xác tương đương SVM hạt nhân Gauss vì vậy không được xác nhận khi thực nghiệm được lặp lại. Trên dữ liệu tín dụng, ngược lại, PWL-LS-SVM cộng tính thuộc nhóm mô hình tốt nhất, cùng với tổ hợp cây tăng cường và thẻ điểm logistic (\cref{tab:credit}).

\emph{(CH2) Chi phí.} Lợi thế về chi phí được xác nhận rõ ràng. Với ánh xạ cộng tính, mô hình nhỏ hơn SVM hạt nhân Gauss từ $2{,}5$ đến $36$ lần trên mười bộ dữ liệu chuẩn và khoảng 690 lần trên dữ liệu tín dụng; thời gian dự đoán nhanh hơn khoảng 200 lần trên Magic và 430 lần trên dữ liệu tín dụng (\cref{tab:cost-measured}). PWL-LS-SVM còn là mô hình phi tuyến huấn luyện nhanh nhất.

\emph{(CH3) Cấu hình.} Ánh xạ cộng tính là lựa chọn mặc định tốt cho dữ liệu dạng bảng: rẻ, dễ diễn giải và thường chính xác nhất trong các dạng ánh xạ. Ánh xạ HH chỉ có lợi khi tương tác giữa các biến đáng kể (bàn cờ, xoắn ốc, Ionosphere); kích thước của nó tăng gần như bậc hai theo số chiều, và quy tắc sinh tham số của bài báo gốc có thể thất bại trên dữ liệu thưa. Độ chính xác bão hòa quanh $M/n=5$--$10$. Hai điều chỉnh nhỏ của đề án --- nút theo phân vị và chuẩn hóa $\|p\|_2=1$ --- cải thiện kết quả và độ ổn định số mà không thay đổi bản chất của phương pháp.

\emph{(CH4) Ứng dụng.} Trên bài toán vỡ nợ thẻ tín dụng, PWL-LS-SVM cộng tính với nút phân vị đạt AUC 0{,}776, chỉ kém một chút so với tổ hợp cây tăng cường (0{,}779) và cao hơn thẻ điểm logistic truyền thống (0{,}773) ở bốn trong năm lần chia, trong khi mô hình chỉ gồm 349 số thực, dự đoán mỗi hồ sơ trong khoảng 0{,}56 micro giây và đọc được như một thẻ điểm tuyến tính từng phần (\cref{fig:credit-shapes}). Khi thêm ràng buộc đơn điệu để mô hình nhất quán với tri thức nghiệp vụ, AUC hầu như không đổi (0{,}775 so với 0{,}776) (Mục~\ref{sec:monotone}).

\paragraph{Khi nào nên dùng PWL-SVM.} Từ các kết quả trên, PWL-SVM cộng tính phù hợp với dữ liệu dạng bảng có số chiều vừa phải, khi mô hình cần nhỏ, dự đoán nhanh hoặc phải giải thích được: chấm điểm tín dụng, các hệ thống nhúng, hay bất cứ bối cảnh nào mà hồi quy logistic đang được dùng vì tính minh bạch. Khi độ chính xác là tiêu chí duy nhất và tài nguyên không bị hạn chế, SVM hạt nhân Gauss, mạng nơ-ron hay tổ hợp cây thường là lựa chọn tốt hơn. Với dữ liệu nhiều chiều, thưa hoặc có tương tác phức tạp, PWL-SVM với các siêu phẳng sinh ngẫu nhiên không phải là công cụ thích hợp; khi đó nên học cả các siêu phẳng, tức là chuyển sang mạng nơ-ron ReLU.

\paragraph{Giới hạn của thực nghiệm.} Kết luận của chương dựa trên mười bộ dữ liệu chuẩn và một bài toán ứng dụng, nên lực của các kiểm định thống kê còn hạn chế. Các lưới tham số được chọn theo khuyến nghị thông dụng và có thể chưa tối ưu cho mọi mô hình; SVM hạt nhân Gauss và mạng nơ-ron trên dữ liệu tín dụng được tinh chỉnh trên mẫu con; thẻ điểm đối chứng dùng một cách rời rạc hóa đơn giản (mười khoảng theo thập phân vị), trong khi các thẻ điểm thực tế thường được xây dựng với các khoảng chọn cẩn thận hơn. Bộ giải LIBLINEAR chính quy hóa nhẹ hệ số chặn, khác với bài toán~\eqref{eq:pwl-c-svm}. Dữ liệu tín dụng được thu thập ở Đài Loan năm 2005, nên các hàm đóng góp tìm được minh họa phương pháp chứ không nhất thiết phản ánh hành vi của khách hàng Việt Nam hiện nay. Cuối cùng, thời gian đo được phụ thuộc vào phần cứng và cách cài đặt.

% ----------------------------------------------------------------------------
\section*{Tiểu kết Chương 3}
\phantomsection\addcontentsline{toc}{section}{Tiểu kết Chương 3}

Chương 3 đã cài đặt PWL-SVM, kiểm chứng tính đúng đắn của cài đặt bằng các phép thử số, và đánh giá phương pháp qua ba lớp thực nghiệm: dữ liệu hai chiều, mười bộ dữ liệu chuẩn và một bài toán ứng dụng. Kết quả cho một bức tranh nhất quán: PWL-SVM không chính xác hơn các mô hình phi tuyến mạnh như SVM hạt nhân Gauss, nhưng nó cho những mô hình nhỏ hơn từ vài lần đến hàng trăm lần, dự đoán nhanh hơn nhiều, huấn luyện bằng một bài toán lồi và, với ánh xạ cộng tính, có thể đọc được từng thành phần. Trên bài toán dự báo vỡ nợ thẻ tín dụng, nơi các tính chất đó được coi trọng, PWL-LS-SVM cộng tính với nút phân vị đạt độ chính xác ngang những mô hình tốt nhất, và với ràng buộc đơn điệu, cho một mô hình nhất quán với tri thức nghiệp vụ mà không mất độ chính xác.
